% Proved after complete root and monodromy audits; exact candidate retained.
\subsection{High quotients in the remaining residual-wild comparison}

Let \(K\) be complete discretely valued of mixed characteristic
\((0,p)\), with algebraically closed residue field, \(p\ge5\),
\(\mu_p\subset K\), and \(e=e_K=p(p-1)\). Put
\(q=p^2\), \(N=q-1\), \(s=(p+1)e=pN\), and let
\(L/K\) be the common tame good extension of degree \(N\).
We compare actual generic \(p^3\)-torsion of two whole potentially
supersingular height-two groups. Assume their residual wild action
is nontrivial. Its unique fixed line is their common canonical
high line \(H\), by Lemma~\ref{lem:B-RD82-Newton}. Write
\[
 0<B_i<pe,\qquad \alpha_i=s-B_i,\qquad \beta_i=B_i/p
                                                        \tag{B-RD91.1}
\]
for the first-level parameters and their high and low coordinate
valuations over \(L\). The common first-level division field
\(R/L\) is cyclic of degree \(p\).

\begin{lemma}[The first-level strict comparison]
\label{lem:B-RD91-first-break}
\status{Proved; full independent root and companion audits passed.}
If \(p\nmid B_1\), then the actual lower break of \(R/L\) is
\[
              b(R/L)=ps-(p+1)B_1.                    \tag{B-RD91.2}
\]
For any second model with the same actual first-level module and
identified high line, one has \(B_1\ge B_2\). In particular,
if both parameters are prime to \(p\), they are equal; if
\(p\mid B_2\), the inequality is strict.
\end{lemma}
\begin{proof}
Choose corresponding nonzero first-level points outside \(H\).
Their stabilizers are the same, and every such point generates
\(R\), since each nonidentity residual unipotent translates it
by a nonzero point of \(H\). The normalized \(R\)-valuations
of their formal coordinates are \(B_i\). Every nonidentity
\(g\in\operatorname{Gal}(R/L)\) displaces the first coordinate
by valuation \(p\alpha_1\). The exact prime-valuation equality
of Lemma~\ref{lem:B-RD60-break} therefore gives (B-RD91.2).
The same lemma's unconditional inequality in the other coordinate
gives
\[
 p(s-B_2)\ge B_2+ps-(p+1)B_1,
 \qquad (p+1)(B_1-B_2)\ge0.                          \tag{B-RD91.3}
\]
If both parameters are prime to \(p\), interchange the models.
If only the first is prime, equality is impossible by divisibility.
These are actual coordinates of the same degree-\(p\) point field;
no bound on the full division field is used.
\end{proof}

\begin{proposition}[Both parameters must be divisible by \(p\)]
\label{prop:B-RD91-divisibility}
\status{Proved; full independent root and companion audits passed.}
If the finite special characters of the two groups differ, then
\(p\mid B_1\) and \(p\mid B_2\).
\end{proposition}
\begin{proof}
Suppose \(p\nmid B_1\). If the second parameter is also prime,
the lemma gives \(B_1=B_2\); equality of the actual high-line
characters \(\zeta^{c_i-B_i}\) then gives \(c_1=c_2\pmod N\).
Thus assume \(p\mid B_2\). The lemma gives \(B_1>B_2\).

Form the actual good integral quotients \(Q_i=M_i/H\), using
the schematic closures of the common line. The product calculation
of Lemma~\ref{lem:B-RD69-quotients} gives their parameters
\[
 D_i=\begin{cases}
       pB_i,&B_i<e,\\
       s-B_i,&B_i>e.
      \end{cases}                                    \tag{B-RD91.4}
\]
At \(B_i=e\) the quotient polygon has one slope, with every
nonzero first-level point at valuation \(e\). The quotient and
the actual semilinear descent are the same constructions as in
B69 and B83: no good integral model over \(K\) is asserted.
The comparison induces an actual \(p^2\)-torsion isomorphism,
because
\[
 Q_i[p^2]=\{P\in M_i[p^3]:p^2P\in H\}/H.           \tag{B-RD91.5}
\]
It also identifies the intrinsic images \(M_i[p]/H\) in
\(Q_i[p]\). Isogeny preserves the unordered finite special
character, so a recovery for these quotients recovers the original
character.

If \(B_i<e\), the intrinsic old image has valuation \(B_i\),
which is the low value \(D_i/p\), while the canonical high value
is \(s-pB_i>e\). These are two distinct stable lines. Over \(L\)
their characters are trivial: the residual image is a \(p\)-group,
whereas every line character is \(\mathbf F_p^\times\)-valued.
The \(p\)-group assertion follows directly because the quotient
torsion is a subquotient of the original \(p^3\)-torsion and the
original residual wild group has order \(p\). Thus \(Q_i[p]\)
is \(L\)-rational. At \(B_i=e\), the one-slope leading-coordinate
map is injective and wild inertia is trivial; the same \(p\)-group
observation again makes \(Q_i[p]\) \(L\)-rational.

If \(B_1<e\), both parameters are below \(e\). Both quotient
modules are therefore \(L\)-rational, and their parameters are
\(D_i=pB_i\). Proposition~\ref{prop:B-RD84-tame}, applied to
their actual \(p^2\)-comparison with \(p\nmid B_1\), recovers
the special characters, a contradiction.

Suppose next \(B_2\le e<B_1\). The second quotient has rational
first-level torsion. The first quotient has parameter
\(D_1=s-B_1\) prime to \(p\); its low points have nonintegral
valuation \(D_1/p\), so its residual division field over \(L\)
has nontrivial wild degree \(p\). This contradicts the actual
quotient first-level isomorphism. This includes the boundary
\(B_2=e\); the prime parameter \(B_1\) cannot equal \(e\).

Finally, if \(e<B_2<B_1\), both quotients have parameters
\(D_i=s-B_i\), with \(D_1\) prime and \(D_2\) divisible by
\(p\). Their identified old images are the canonical high lines,
of valuation \(B_i=s-D_i\). The first quotient has nontrivial
residual wild action, hence so does the second. The first-level
lemma applied to this actual quotient comparison gives
\(D_1>D_2\), opposite to \(B_1>B_2\). All cases are excluded.
\end{proof}

\begin{proposition}[The crossing parameters at \(p=5\)]
\label{prop:B-RD91-parameters}
\status{Proved; full independent root and companion audits passed.}
For \(p=5,e=20\), normalize the nonsquare common finite determinant
as in B77. A remaining different-character residual-wild comparison
has
\[
 B_i=5b_i,\qquad b_i\in\{1,3,5,7,9,11,13,15,17,19\},
 \qquad b_1\equiv b_2\pmod4.                         \tag{B-RD91.6}
\]
The two parameters cannot both be below \(20\). If they straddle
\(20\), their ordered values, with \(B_1<20<B_2\), are exactly
among the two necessary rows
\[
                 (B_1,B_2)=(5,45),\ (15,95).          \tag{B-RD91.7}
\]
Their common high quotients have rational first-level torsion and
exchanged canonical high lines, with unordered parameters
\(\{25,75\}\). If both original parameters exceed \(20\), their
common high quotients have nontrivial residual wild action.
\end{proposition}
\begin{proof}
The preceding proposition gives \(B_i=5b_i\), with
\(1\le b_i\le19\). Normalize \(c_i\equiv1\pmod4\).
The high-line character is \(\mathbf F_5^\times\)-valued, so
\(c_i-B_i\equiv0\pmod6\). Consequently \(B_i\), and hence
\(b_i\), is odd. Equality of the high-line characters and finite
determinants gives \(B_1\equiv B_2\pmod4\), hence the displayed
congruence. Equal parameters recover the finite character, so they
are distinct here. The only parameters below \(20\) are \(5,15\),
whose quotients by five are not congruent modulo four. The boundary
\(B_i=20\) is excluded by oddness.

In a crossing row the smaller quotient has rational first-level
torsion, so the other quotient does as well. Their parameters are
\(25b_1\) and \(5(24-b_2)\). Proposition B84 forces
\(5\mid(24-b_2)\), since otherwise one parameter divided by five
is prime to five and the quotient comparison recovers the character.
The odd possibilities \(5\le b_2\le19\) are \(9,19\). The
congruence modulo four then gives \((b_1,b_2)=(1,9),(3,19)\),
which yields (B-RD91.7). The old first-level image is low in the
smaller original model's quotient and high in the other, proving
the stated exchange and the actual quotient parameters.

If both original parameters exceed \(20\), their quotient
parameters are \(D_i=5(24-b_i)\), and their identified old images
are the common quotient high line. Suppose the quotient residual
action were tame, hence trivial over \(L\). B84 would force
both \(24-b_i\) divisible by five. Thus both \(b_i\) would lie
in \(\{9,19\}\), and their congruence modulo four would force
equality. This recovers the character, a contradiction. Therefore
the quotient residual wild action is nontrivial.
\end{proof}

\begin{theorem}[Only the two largest layer shapes remain above the index]
\label{thm:B-RD91-large}
\status{Proved; full independent root and companion audits passed.}
At \(p=5,e=20\), assume additionally that both \(B_i>20\).
Let \(W\) be the common actual \(125\)-division image over \(L\),
and let \(W_1,W_2\subset M_2(\mathbf F_5)\) be its first and
terminal congruence layers. In a basis with the common high line
\(\mathbf F_5e_1\),
\[
 \mathfrak{sl}_2(\mathbf F_5)\subset W_1\subset W_2,
 \qquad (\dim W_1,\dim W_2)\in\{(3,4),(4,4)\}.
                                                        \tag{B-RD91.8}
\]
Thus its full wild order is respectively \(5^8\) or \(5^9\).
\end{theorem}
\begin{proof}
The original residual image consists of matrices
\(1+uE_{12}\). For \(g\in W\) put
\[
 u(g)=\rho_{12}(g)\bmod5,\qquad
 v(g)=\rho_{21}(g)/5\bmod5.                          \tag{B-RD91.9}
\]
Both are additive homomorphisms \(W\to\mathbf F_5\): the
diagonal entries are one modulo five and the lower entry is
divisible by five. In the actual high-quotient Tate lattice,
with basis \((e_1/5,e_2)\), the residual lower entry is \(v(g)\),
as in (B-RD83.6). The preceding proposition says that \(v\ne0\).

Choose a prime-to-five order-24 lift of a tame generator in the
actual finite division group, and diagonalize its action on the
\(125\)-torsion by the coprime factors of \(X^{24}-1\), exactly
as in B90. Let its residual eigenvalues be \(\alpha,\beta\).
Its ratio \(r=\alpha/\beta\) is nonsquare in \(\mathbf F_5\),
so \(1,r,r^{-1}\) are distinct. Conjugation gives weights
\(r,r^{-1}\) on \(u,v\). If \(v\) vanished on
\(W^1=\ker u\), it would factor as a nonzero scalar multiple
of \(u\), contradicting those distinct weights. Hence \(W_1\)
has a nonzero lower off-diagonal entry. The tame spectral
projector gives \(E_{21}\in W_1\). Conjugating it by the
residual matrices \(1+tE_{12}\), and using \(t=1,-1\), gives
\(E_{12}\) and \(\operatorname{diag}(1,-1)\) as well.
Thus \(\mathfrak{sl}_2\subset W_1\). Actual fifth powers give
\(W_1\subset W_2\), and
\[
                  |W|=5^{1+\dim W_1+\dim W_2}.        \tag{B-RD91.10}
\]

It remains to exclude \(W_1=W_2=\mathfrak{sl}_2\). In that
case determinants on \(W^1\) are one modulo 25. The map
\(g\mapsto(\det g-1)/25\bmod5\) is additive, vanishes on the
terminal subgroup \(1+25W_2\), and therefore factors through a
linear functional \(\ell:W_1\to\mathbf F_5\). It is invariant
under both tame conjugation and residual unipotent conjugation.
The two nontrivial tame weights kill \(\ell(E_{12})\) and
\(\ell(E_{21})\); conjugation of \(E_{21}\) by \(1+E_{12}\)
then kills \(\ell(\operatorname{diag}(1,-1))\). Thus
\(\det(W^1)=1\pmod{125}\).

The determinant on \(W\) is the actual cyclotomic character,
since its finite tame factor is trivial on \(L\). Its image
therefore has order at most five. The exact cyclotomic tower
argument in B86 shows that it has order five, that
\(\mu_{25}\subset K\), and that, for a Cohen coefficient field
\(K_0\subset K\), one has \(K=K_0(\mu_{25})\). Indeed the
absolute index is exactly 20, while \(\mu_{125}\) requires
index 100; a tame degree-24 extension cannot change the remaining
five-primary cyclotomic degree. Since the determinant is trivial
on \(W^1\), its degree-five field is precisely
\(R=F^{W^1}\), the original residual division field.

The exact cyclotomic calculation of B86 now gives
\(b(R/L)=24^2=576\). On the other hand, take an actual low
first-level point. Its normalized \(R\)-coordinate has value
\(B_i\), and every nonidentity residual automorphism displaces
it by value \(5(s-B_i)\). The unconditional bound of B60 gives
\[
             b(R/L)\le5s-6B_i=600-6B_i<576,
                        \qquad B_i>20,                \tag{B-RD91.11}
\]
a contradiction. This excludes the only layer shape of total
order \(5^7\), and proves (B-RD91.8).
\end{proof}

\begin{remark}[Exact remaining scope]
\label{rem:B-RD91-scope}
\status{Proved; full independent root and companion audits passed.}
The crossing rows (B-RD91.7) are necessary patterns, not examples.
The theorem addresses only the both-large residual-wild branch;
it neither excludes nor constructs its two remaining layer shapes.
The irreducible residual branch is unchanged. The unrestricted
first-wild least index at \(p=5,e=20\) remains in \(\{2,3\}\).
All quotient comparisons have the actual levels displayed in
(B-RD91.5); in particular a quotient \(25\)-comparison is never
silently treated as an actual \(125\)-comparison.

The proof rereads the complete internal B60 fixed-coefficient
lemma, B69 quotient formula, B82 residual normalization, B84
point-field recovery, and B86 cyclotomic-field calculation.
No fresh original-source reading is claimed here. The retained
standalone program
\path{code/reduction_descent/verify_residual_wild_high_quotient.py}
reconstructs 78,761 inverse-Newton quotient distributions, including
all nine tested equality boundaries, and checks 210,956 complete
prime/divisible case rows over the four stated small primes.
Its independent finite-character enumeration retains 32 normalized
distinct-character rows, the four ordered crossing rows representing
(B-RD91.7), and the 20 both-large rows. It reconstructs all 250
linear functionals on the two tame forms of the trace-zero plane,
excludes all eight nonzero intertwiners between the opposite tame
characters, and checks 75,571 cyclotomic break inequalities.
The first author run passed at 03:02:51.567982 UTC on 20 November
2026 in 0.6828030579927145 seconds. No executed computation failed.
The program, full output, and pre-certificate source are retained.
These finite checks do not replace the uniform field and quotient
proofs, and no parameter row is treated as a formal-group existence
assertion.

\smallskip
\noindent\emph{Independent audits.}
The complete root and monodromy proof and exact-code audits passed.
Their literal author-output-preserving reruns took
\(0.69858235999709\) and
\(0.6949913400021615\) seconds, respectively.
Both reread the complete invoked B60, B69, B82, B84 and B86
dependency proofs. No additional original-source opening is claimed.
The exact audited candidate, pre-promotion copy, programs, outputs,
two full audit records and promotion diff remain retained.
\end{remark}
